\section[Cote 5 : l'algèbre artinienne universelle à puissances divisées]{Cote 5 : l'algèbre artinienne universelle à puissances divisées\protect\verifie{Folder5}}\label{sec:5}

La cote 5, \emph{Cristaux et cohomologie de De Rham : généralités, correspondances Deligne et Berthelot}, que l'inventaire date de 1965--1969, porte sur les puissances divisées et les sites cristallins. À la page~4, l'algèbre locale artinienne universelle sur $W(k)$ dont l'idéal maximal porte des puissances divisées nulles en degrés $\geq p^{\nu+1}$ est présentée comme $W[\Lambda_0, \dots, \Lambda_\nu]/(\Lambda_i^p - q_i\Lambda_{i+1}, \Lambda_\nu^p, \Lambda_0 - p)$, où $\Lambda_i = \gamma_{p^i}(p)$ et $q_i = (p^{i+1})!/((p^i)!)^p$ (\lot{5}{1}{4}) ; la relation est $\gamma_{p^i}(x)^p = q_i\, \gamma_{p^{i+1}}(x)$. Une note de la lecture s'appuie sur $v_p(q_\nu) = 1$.

\begin{theoreme}[Énoncé de l'édition, d'après la page 4\piste{5-universal-artinian-pd-normal-form}]
Pour tout entier $p$ et tout $i$, $((p^i)!)^p\, q_i = (p^{i+1})!$ : $q_i$ est un entier, le coefficient multinomial de $p$ blocs de taille $p^i$. Si $p$ est premier, $(p-1)\, v_p((p^i)!) = p^i - 1$ et $v_p(q_i) = 1$.
\end{theoreme}

\noindent L'énoncé tient. La valuation vient de la formule de Legendre sous la forme $v_p((pm)!) = v_p(m!) + m$ ; voir les théorèmes \texttt{factorial\_eq\_mul\_q}, \texttt{q\_eq\_multinomial} et \texttt{padicValNat\_q} de \href{\leanblob/Folder5.lean}{\texttt{Folder5.\allowbreak{}lean}}.

\noindent\emph{L'universalité (page~4).} La page place l'énoncé d'universalité entre crochets, sans démonstration ; la lecture en esquisse une par l'enveloppe à puissances divisées. La vérification en donne une autre, qui n'utilise que mathlib. On travaille sur un anneau intègre $W$ de caractéristique nulle où les entiers premiers à $p$ sont inversibles : c'est le cas de $W(k)$ et de $\mathbf Z_p$ (\texttt{padicInt\_isUnit\_natCast}). On pose
\[
A^{(\nu)} = W[\Lambda_0, \ldots, \Lambda_\nu] \big/ \bigl(\Lambda_i^{p} - q_i\Lambda_{i+1}\ (i < \nu),\ \Lambda_\nu^{p}\bigr), \qquad W^{(\nu)} = A^{(\nu)}/(\Lambda_0 - p),
\]
qui est bien le quotient de $W[\Lambda]$ par toutes les relations de la page (\texttt{presentation}), et $J = (\Lambda_0, \ldots, \Lambda_\nu) \subset W^{(\nu)}$.

\begin{theoreme}[Grothendieck, page 4, universalité de la présentation\piste{5-universal-artinian-pd-normal-form}]
L'idéal $J$ de $W^{(\nu)}$ porte une structure de puissances divisées $\delta$ telle que $\delta_{p^i}(p) = \Lambda_i$ pour $0 \leq i \leq \nu$ et $\delta_n(y) = 0$ pour $y \in J$ et $n \geq p^{\nu+1}$. Pour toute $W$-algèbre $T$ munie d'un idéal $I$ à puissances divisées $\gamma$ tel que $p \in I$ et $\gamma_n(p) = 0$ pour $n \geq p^{\nu+1}$, il existe un unique morphisme de $W$-algèbres $W^{(\nu)} \to T$ compatible aux puissances divisées ; il envoie $\Lambda_i$ sur $\gamma_{p^i}(p)$.
\end{theoreme}

\begin{proof}
\emph{Un plongement.} Soit $K$ le corps des fractions de $W$ et $B = K[X]/(X^{p^{\nu+1}})$. L'application $\psi : \Lambda_i \mapsto X^{p^i}/(p^i)!$ respecte les relations, puisque $((p^i)!)^p q_i = (p^{i+1})!$, et définit $\psi : A^{(\nu)} \to B$. Elle est injective. D'une part les monômes $\Lambda^a$ à exposants $0 \leq a_i < p$ engendrent $A^{(\nu)}$ comme $W$-module : leur $W$-module est stable par multiplication par chaque $\Lambda_j$, grâce à la retenue de la page~8 ($\Lambda_j^{p} = q_j\Lambda_{j+1}$, ou $0$ si $j = \nu$). D'autre part $\psi(\Lambda^a)$ est un multiple non nul de $X^{\sum a_i p^i}$, et ces exposants sont distincts et $< p^{\nu+1}$ (écriture en base $p$). Un élément $\sum w_a \Lambda^a$ de noyau nul a donc tous ses $w_a$ nuls.

\emph{Les puissances divisées de $A^{(\nu)}$.} Comme $B$ est une $\mathbf Q$-algèbre, tout idéal de $B$ porte les puissances divisées $x^n/n!$. On les ramène sur $\bar I = (\Lambda_0, \ldots, \Lambda_\nu) \subset A^{(\nu)}$ par $\psi$ : il suffit que $\gamma_n(\psi(y))$ soit dans $\psi(\bar I)$ pour $y \in \bar I$. Par les formules de la somme et du produit, il suffit de le voir pour $y = \Lambda_i$, et comme $\gamma_n(\gamma_{p^i}(X))$ est un multiple entier de $\gamma_{np^i}(X)$, il suffit que chaque $\gamma_k(X) = X^k/k!$ vienne de $A^{(\nu)}$. C'est une récurrence sur $k$ : si $p^e$ est la plus grande puissance de $p$ divisant $k$, on a $\binom{k}{p^e}\gamma_k = \gamma_{p^e}\,\gamma_{k-p^e}$, et $\binom{k}{p^e} \equiv k/p^e \not\equiv 0 \pmod p$ par le théorème de Lucas, donc inversible dans $W$ ; et $\gamma_{p^e}(X) = \psi(\Lambda_e)$ si $e \leq \nu$, tandis que $\gamma_k(X) = 0$ si $k \geq p^{\nu+1}$.

\emph{Passage au quotient.} L'augmentation $\varepsilon : A^{(\nu)} \to W$, $\Lambda_i \mapsto 0$, a pour noyau $\bar I$. Si $(\Lambda_0 - p)a \in \bar I$, alors $-p\,\varepsilon(a) = 0$, donc $a \in \bar I$ ; et $\gamma_n((\Lambda_0 - p)a) = (\Lambda_0 - p)^n\gamma_n(a)$. Ainsi $(\Lambda_0 - p) \cap \bar I$ est un sous-idéal à puissances divisées, et celles-ci passent à $J = \bar I\,W^{(\nu)}$. On a $\delta_{p^i}(p) = \delta_{p^i}(\Lambda_0) = \Lambda_i$ par injectivité de $\psi$, et $\delta_n(y) = 0$ pour $n \geq p^{\nu+1}$ parce que $\psi(y)$ est un multiple de $X$.

\emph{L'universalité.} Le morphisme $\Lambda_i \mapsto \gamma_{p^i}(p)$ respecte les relations : $\gamma_{p^i}(p)^p = q_i\,\gamma_{p^{i+1}}(p)$ (\texttt{dpow\_pow\_eq\_q}), $\gamma_{p^\nu}(p)^p = q_\nu\gamma_{p^{\nu+1}}(p) = 0$ et $\gamma_1(p) = p$. Il commute aux $\delta_k(p)$ par la même récurrence sur $k$, et donc aux puissances divisées des générateurs $\Lambda_i$, ce qui suffit. Un morphisme compatible envoie nécessairement $\Lambda_i = \delta_{p^i}(p)$ sur $\gamma_{p^i}(p)$, d'où l'unicité.
\end{proof}

\noindent L'énoncé tient, sous une forme plus forte que celle de la page : ni la locali\-té, ni l'artinianité, ni le corps résiduel n'interviennent, et l'hypothèse sur $T$ ne porte que sur les $\gamma_n(p)$. Lean : \texttt{universal}, \texttt{dpWν}, \texttt{dpWν\_p}, \texttt{dpWν\_eq\_zero}, \texttt{presentation} ; pour les étapes, \texttt{boxSpan\_eq\_top} (la retenue), \texttt{deg\_injective}, \texttt{psi\_injective}, \texttt{dpow\_split} et \texttt{choose\_pow\_mul\_pow\_modEq} (Lucas), \texttt{psi\_hmem}, \texttt{dpIbar}, \texttt{isSubDPIdeal\_Kp}, dans \href{\leanblob/Folder5.lean}{\texttt{Folder5.\allowbreak{}lean}}.

\begin{theoreme}[Grothendieck, pages 5 et 6, la suffisance du lemme\piste{5-universal-artinian-pd-normal-form}]
Soient $\lambda_0, \ldots, \lambda_\nu \in W$ avec $\lambda_0 = p$ et $\lambda_i^p = q_i\lambda_{i+1}$ (ainsi $\lambda_i = p^{p^i}/(p^i)!$), et $\Theta_i = \Lambda_i - \lambda_i$. Dans $W^{(\nu)}$, on a $p^i\Theta_i = 0$ pour $0 \leq i \leq \nu$ et $p^\nu\lambda_\nu^p = 0$ ; si $\lambda_\nu = p^m u$ avec $u$ inversible, $p^{\nu + pm} = 0$.
\end{theoreme}

\begin{proof}
Récurrence sur $i$, comme à la page~6 : $\Theta_0 = \Lambda_0 - p = 0$. Si $p^i\Theta_i = 0$, la différence $\Lambda_i^p - \lambda_i^p = q_i\Theta_{i+1}$ est un multiple de $\Lambda_i - \lambda_i = \Theta_i$, donc $p^i q_i \Theta_{i+1} = 0$ ; comme $q_i = p\,r_i$ avec $r_i$ premier à $p$ ($v_p(q_i) = 1$), $p^{i+1}\Theta_{i+1} = 0$. Pour $i = \nu$, $\Lambda_\nu^p = 0$ donne de même $p^\nu\lambda_\nu^p = 0$. On n'a pas besoin de développer le binôme.
\end{proof}

\noindent Pour $\lambda_\nu = p^{p^\nu}/(p^\nu)!$, $m = p^\nu - (p^\nu - 1)/(p-1)$ (par \texttt{sub\_one\_mul\_padicValNat\_factorial\_pow}) et $\nu + pm = c(\nu)$ : c'est le \q{$p^{c(\nu)} \in I$} de la page~6. Lean : \texttt{p\_pow\_mul\_Θ\_mem}, \texttt{p\_pow\_mul\_lam\_pow\_mem}, \texttt{p\_pow\_mem\_Kp}, énoncés comme appartenances à l'idéal $(\Lambda_0 - p)$ de $A^{(\nu)}$.

La nécessité, que la page laisse \q{au lecteur} avec un point d'exclamation en marge, est démontrée elle aussi. Pour un exposant $a = (a_0, \ldots, a_\nu)$, $0 \leq a_i < p$, notons $\operatorname{ord}(a)$ l'indice du premier exposant non nul de $a$, et $\operatorname{ord}(0) = c(\nu) = \nu + pm$.

\begin{theoreme}[Grothendieck, pages 5 et 6, le lemme et son corollaire\piste{5-universal-artinian-pd-normal-form}]
Soient $\lambda_i$ comme ci-dessus, avec $\lambda_\nu = p^m u$, $u$ inversible, et $\Theta_i = \Lambda_i - \lambda_i$.
\begin{enumerate}
\item Les monômes $\Theta^a$, $0 \leq a_i < p$, forment une base du $W$-module $A^{(\nu)}$.
\item \emph{(Le lemme de la page 5.)} $\sum_a w_a\Theta^a$ appartient à $(\Lambda_0 - p)$ si et seulement si $p^{\operatorname{ord}(a)}$ divise $w_a$ pour tout $a$ : $w_a \in p^iW$ si $a_i$ est le premier exposant non nul, et $w_0 \in p^{c(\nu)}W$.
\item \emph{(Le corollaire de la page 6.)} $W^{(\nu)} \simeq \bigoplus_a W/p^{\operatorname{ord}(a)}W$ comme $W$-module.
\item Si $p$ n'est pas inversible dans $W$, $p^k = 0$ dans $W^{(\nu)}$ si et seulement si $k \geq c(\nu)$ : l'ordre de $1$ est exactement $p^{c(\nu)}$.
\item Si $W$ est un anneau de valuation discrète d'uniformisante $p$, $W^{(\nu)}$ est de longueur $\sum_a \operatorname{ord}(a) = p + p^2 + \cdots + p^\nu + pm$.
\end{enumerate}
\end{theoreme}

\begin{proof}
\emph{La retenue dans les $\Theta_i$.} En développant $\Lambda_i^p = (\Theta_i + \lambda_i)^p$ et en utilisant $\lambda_i^p = q_i\lambda_{i+1}$, on trouve
\[
\Theta_i^p = q_i\Theta_{i+1} - \sum_{0<k<p}\binom pk\lambda_i^{p-k}\Theta_i^k \quad (i < \nu), \qquad \Theta_\nu^p = -\lambda_\nu^p - \sum_{0<k<p}\binom pk\lambda_\nu^{p-k}\Theta_\nu^k .
\]
Pour $i < \nu$ le terme constant $\lambda_i^p$ disparaît : c'est ce qui rend la retenue compatible avec les conditions a) et b).

\emph{Un idéal.} Soit $S_n$ le sous-$W$-module engendré par les $p^{\min(\operatorname{ord}(a), n)}\Theta^a$. On montre, par récurrence descendante sur $j$, que $p^{\min(\operatorname{ord}(a), j, n)}\,\Theta_j\Theta^a \in S_n$ pour tout $a$. Si $a_j < p - 1$, $\Theta_j\Theta^a$ est un monôme de la base, d'ordre au plus $\min(\operatorname{ord}(a), j)$. Si $a_j = p - 1$, soit $a'$ l'exposant $a$ où l'on a remplacé $a_j$ par $0$ ; alors $\Theta_j\Theta^a = \Theta_j^p\Theta^{a'}$, et l'on remplace $\Theta_j^p$ par la formule ci-dessus. Les termes en $\Theta_j^k\Theta^{a'}$, $0 < k < p$, sont des monômes de la base d'ordre au plus $\operatorname{ord}(a)$. Le terme $q_j\Theta_{j+1}\Theta^{a'}$ relève de l'hypothèse de récurrence : l'ordre ne monte que de $1$ au plus, de $j$ à $j+1$, et le facteur $p$ de $q_j$ le compense. Pour $j = \nu$, le terme $\lambda_\nu^p\Theta^{a'}$ porte le facteur $p^{pm}$, qui compense le passage de l'ordre $\nu$ à l'ordre $c(\nu) = \nu + pm$ quand $a' = 0$. Comme $\Lambda_j = \Theta_j + \lambda_j$, $S_n$ est stable par multiplication par les $\Lambda_j$ : c'est un idéal.

\emph{La base.} $S_0$ est un idéal qui contient $1$ : les $\Theta^a$ engendrent $A^{(\nu)}$. Comme les $\Lambda^a$ en forment une base (c'est la seconde moitié de l'injectivité de $\psi$), l'application $(w_a) \mapsto \sum_a w_a\Theta^a$ est une surjection d'un module libre de rang $p^{\nu+1}$ sur un module libre de même rang, donc une bijection : sur un anneau commutatif, un endomorphisme surjectif d'un module de type fini est injectif.

\emph{Le lemme et ses conséquences.} $S_{c(\nu)}$ est un idéal qui contient $\Theta_0 = \Lambda_0 - p$, donc contient $(\Lambda_0 - p)$ : c'est la nécessité. La suffisance (théorème précédent) donne l'inclusion inverse. L'unicité des coordonnées dans la base des $\Theta^a$ donne alors (2), puis (3). Pour (4), on applique (2) à $p^k = p^k\,\Theta^0$ : il faut et il suffit que $p^{c(\nu)}$ divise $p^k$. Pour (5), $W/p^kW$ est de longueur $k$, et l'on compte : si $a_0 \neq 0$, l'ordre est $0$, et si $a_0 = 0$, l'ordre est celui de $(a_1, \ldots, a_\nu)$ augmenté de $1$. Par récurrence sur $\nu$, $\sum_a \operatorname{ord}(a) + \nu + 1 = 1 + p + \cdots + p^\nu + c(\nu)$.
\end{proof}

\noindent L'énoncé tient. La seule propriété des $q_i$ dont la nécessité ait besoin est $p \mid q_i$. La longueur $p + \cdots + p^\nu + pm$ vaut, par télescopage, la formule $c(\nu) + \sum_{i=1}^{\nu} i(p-1)p^{\nu-i}$ de la lecture ; pour $p = 2$ et $\nu = 1$, on retrouve la longueur $4$ et $c(1) = 3$ de la note de la lecture. Pour les $\lambda_i$ usuels, définis par $\lambda_i\,(p^i)! = p^{p^i}$, on a $m = p^\nu - v_p((p^\nu)!)$, et pour $W = \mathbf Z_p$ l'ordre de $1$ est exactement $p^{c(\nu)}$. Lean : \texttt{normalForm}, \texttt{tmono\_linearIndependent}, \texttt{span\_tmono}, \texttt{WνEquiv}, \texttt{p\_pow\_eq\_zero\_iff}, \texttt{length\_Wν}, \texttt{length\_Wν\_eq} ; pour les étapes, \texttt{Θ\_pow\_castSucc} et \texttt{Θ\_pow\_last} (la retenue), \texttt{Θ\_mul\_tmono\_mem} et \texttt{mul\_mem\_S} (l'idéal), \texttt{mono\_linearIndependent}, \texttt{Kp\_le\_S} (la nécessité), \texttt{S\_le\_Kp}, \texttt{sum\_ord} (le dénombrement) ; pour les $\lambda_i$ usuels, \texttt{exists\_lamStd}, \texttt{p\_pow\_eq\_zero\_iff\_std}, \texttt{length\_Wν\_std} et \texttt{padicInt\_Wν}, dans \href{\leanblob/Folder5.lean}{\texttt{Folder5.\allowbreak{}lean}}.

Reste hors de la vérification le lemme de la page~7 :

\begin{encadre}
\begin{theoreme*}[Grothendieck, page 7, les quotients $W^{(\nu,r)}$\piste{5-universal-artinian-pd-normal-form}]\nonverifie
Pour $1 \leq r \leq \nu$, soit $W^{(\nu,r)}$ le quotient de $W^{(\nu)}$ par $\Theta_1, \ldots, \Theta_r$. Alors $\sum w_a\Theta^a$ est dans $(\Theta_0, \ldots, \Theta_r)$ si et seulement si $w_a \in p^{i-r}W$ lorsque le premier exposant non nul $a_i$ a un indice $i > r$, et $w_0 \in p^{c(\nu)-r}W$.
\end{theoreme*}
\end{encadre}
{\small La forme normale de $W^{(\nu)}$ est démontrée ci-dessus. Pour les quotients $W^{(\nu,r)}$, l'édition n'a fait qu'une vérification à la machine, pour quelques $(p, \nu)$.\par}

\begin{encadre}
\begin{theoreme*}[Grothendieck, pages 49 et 50, dualité étale pour les sites annelés\piste{5-deligne-letter-ringed-duality}]\nonverifie
Soient $f : (X, B) \to (Y, A)$ un morphisme lisse compactifiable de sites annelés, $A$ un faisceau d'anneaux sur $Y$ tué par un entier $n$ premier aux caractéristiques résiduelles, et $B$ un faisceau d'anneaux muni de $f^{-1}(A) \to B$. La dualité étale vaut pour $f$, avec $f^!(K) = R\mathrm{Hom}_{f^{-1}(A)}(B, f^{-1}(K) \otimes T_{X/Y}[2d])$, et se ramène aux cas $f^{-1}(A) \simeq B$ et $f = \mathrm{id}$. La lettre du 10~décembre~1965 l'énonce sans démonstration.
\end{theoreme*}
\end{encadre}
{\small Mathlib n'a pas la dualité étale. On se reportera aux notes de Grothendieck (\lot{5}{3}{49}).\par}
